\documentclass[11pt,a4paper]{article}
\usepackage[margin=2cm]{geometry}
\usepackage{fontspec}
\usepackage[italian]{babel}
\usepackage{xcolor,graphicx,booktabs,longtable,hyperref,fancyhdr}
\setmainfont{NimbusSans-Regular.otf}[Path=sorgenti/font/,BoldFont=NimbusSans-Bold.otf,ItalicFont=NimbusSans-Italic.otf]
\setmonofont{lmmono10-regular.otf}[Path=sorgenti/font/,Scale=.88]
\definecolor{Ink}{HTML}{253344}\definecolor{Navy}{HTML}{15283E}\definecolor{Teal}{HTML}{087C74}\definecolor{Purple}{HTML}{5146A5}
\color{Ink}\hypersetup{colorlinks=true,linkcolor=Teal,urlcolor=Teal,pdftitle={Rilettura delle 861 slide · R6}}
\setlength{\parindent}{0pt}\setlength{\parskip}{8pt}\setlength{\headheight}{15pt}
\pagestyle{fancy}\fancyhf{}\fancyhead[L]{\small Geologia · Statistica e coding R}\fancyhead[R]{\small Rilettura R6}\fancyfoot[L]{\small 9 novembre 2026}\fancyfoot[R]{\thepage}
\newcommand{\titolo}[1]{{\color{Navy}\LARGE\bfseries #1\par}\vspace{4pt}{\color{Teal}\hrule height 1.4pt}\vspace{12pt}}
\begin{document}
\titolo{Esito della rilettura integrale}
{\Large\color{Teal}\textbf{861 slide esaminate · 357 con interventi}}

La rilettura ha considerato ciascuna pagina dei 28 fascicoli, verificando congiuntamente testo, formalismo, esempi e composizione. Le 504 pagine che non richiedevano correzioni conservano il contenuto della R5. I 357 interventi riguardano il contenuto o la rappresentazione; l’aggiornamento dell’etichetta di edizione sulle copertine è escluso dal conteggio.

\subsection*{Autonomia della lettura}
Le spiegazioni sono state integrate laddove un simbolo, un’unità, una convenzione o un’ipotesi restavano impliciti. Sono stati eliminati rinvii editoriali che obbligavano a ricostruire la storia delle versioni. L’autonomia riguarda il ragionamento sviluppato nella pagina e presuppone le conoscenze previste dal modulo: una slide sulla covarianza HAC conserva, per esempio, i prerequisiti di regressione e algebra matriciale.

\subsection*{Intelligibilità delle rappresentazioni}
Formule e diagrammi sono stati confrontati con gli esempi numerici e con il codice pertinente. Il controllo ha riguardato anche orientamento degli assi, significato delle frecce, distinzione fra variabili e parametri, unità di campionamento e corrispondenza fra geometria illustrata e operazione statistica. Gli interventi comprendono correzioni concettuali, oltre a integrazioni esplicative.

\subsection*{Qualità compositiva}
Il risultato conserva l’impostazione di una dispensa visiva: prosa argomentata, gerarchia stabile, figure vettoriali e palette coerente. Sono stati esaminati contrasto, proporzioni, leggibilità delle etichette e rapporto fra spazio scritto e spazio figurato. Il template, le dimensioni delle pagine e i file dei caratteri sono rimasti invariati. Il giudizio estetico è qualitativo e non viene trasformato in un punteggio arbitrario.

\subsection*{Metodo e documentazione}
Il testo di ogni slide è stato riletto e il relativo PDF esaminato nelle tavole di contatto, ricorrendo a ingrandimenti nei passaggi dubbi. Le pagine corrette sono state compilate e ricontrollate. Il registro allegato contiene 861 osservazioni individuali; la verifica geometrica automatica affianca questo esame.

\newpage
\titolo{Copertura e interventi per fascicolo}
La numerazione rimane invariata rispetto alla R5. La colonna «Corrette» conta le pagine sulle quali è intervenuta questa rilettura.
\small
\begin{longtable}{p{1.25cm}p{10.95cm}rr}
\toprule Modulo & Fascicolo & Slide & Corrette\\\midrule\endhead
T00 & Dal codice al calcolo scientifico & 33 & 15\\
M01 & Pensare con i dati & 58 & 17\\
M02 & R e gestione dei dati & 43 & 15\\
M03 & Descrivere e visualizzare i dati & 33 & 12\\
M04 & Probabilità, variabilità e simulazione & 35 & 13\\
M05 & Campionamento, stima e incertezza & 105 & 40\\
M06 & Confrontare gruppi e valutare ipotesi & 35 & 21\\
M07 & Relazioni tra variabili e regressione & 33 & 18\\
M08 & Analisi multivariata dei campioni & 33 & 22\\
M09 & Serie temporali ambientali & 33 & 15\\
M10 & Previsione e progetto statistico & 33 & 18\\
M11 & Statistica spaziale & 32 & 15\\
M12 & Raster, telerilevamento e morfometria & 30 & 12\\
M13 & Cronologie e archivi sedimentari & 30 & 9\\
M14 & Estremi, pericolosità e rischio & 31 & 12\\
M15 & Riproducibilità e progetto di analisi & 27 & 10\\
M16 & Classificazione e validazione & 29 & 16\\
P01 & Geometria dei pori e permeabilità & 21 & 9\\
P02 & Foraminiferi, coppe e campagne & 22 & 7\\
P03 & Il Nilo: livelli e dipendenza temporale & 21 & 8\\
P04 & Un geyser: associazione e previsione & 21 & 7\\
P05 & CO\textsubscript{2}: tendenza e calendario stagionale & 22 & 8\\
P06 & Dai campioni alla validazione spaziale & 17 & 6\\
P07 & Terreno, pendenze e sintesi territoriali & 17 & 7\\
P08 & Collocare nel tempo il massimo di un proxy & 18 & 6\\
P09 & Un livello raro e la sua incertezza & 16 & 8\\
P10 & Un’analisi ricostruibile e verificabile & 16 & 3\\
P11 & Classificare nuove stazioni sintetiche & 17 & 8\\
\midrule\textbf{Totale} & & \textbf{861} & \textbf{357}\\\bottomrule
\end{longtable}
\normalsize
\newpage
\titolo{Un difetto grafico con origine nel codice}
\textbf{M12, pagine 21–22 · Ombreggiatura del terreno}

La figura presentava striature che non appartenevano alla morfologia simulata. La funzione \texttt{pmax}, ricevendo come primo argomento lo scalare zero, perdeva la struttura matriciale; la successiva trasposizione riordinava quindi i valori in modo incoerente con le celle.

La correzione conserva la matrice e verifica, cella per cella, il prodotto fra normale al terreno e direzione della luce per entrambi gli azimut, 315 e 135 gradi. Il laboratorio M12 N12, che conteneva la stessa operazione, è stato corretto e rieseguito con esito positivo.

\begin{minipage}[t]{.48\textwidth}\textbf{\color{Purple}R5 · prima della correzione}\\[6pt]
\includegraphics[width=\linewidth]{audit/evidenze/M12_F06_R5.pdf}
\end{minipage}\hfill
\begin{minipage}[t]{.48\textwidth}\textbf{\color{Teal}R6 · matrice conservata}\\[6pt]
\includegraphics[width=\linewidth]{figure/M12/F06.pdf}
\end{minipage}

L’esempio mostra perché la valutazione estetica di una figura scientifica debba comprendere la correttezza della rappresentazione: una trama artificiale avrebbe potuto essere interpretata come una struttura del terreno.

\textbf{Riscontri conservati.} I sorgenti si trovano in \texttt{R/M12\_figure.R} e \texttt{laboratori/M12/M12\_N12.R}; i resoconti di esecuzione sono in \texttt{verifiche/M12\_figure.log} e \texttt{verifiche/M12\_N12\_R6.log}.

\newpage
\titolo{Leggibilità e autonomia locale}
\textbf{P07, pagina 14 · Numerazione dei blocchi}

I numeri sovrapposti alla carta potevano confondersi con i colori sottostanti. Il fondo chiaro delle etichette rende uniforme il contrasto, mantenendo la palette della carta e la numerazione dei dodici blocchi.

\begin{minipage}[t]{.48\textwidth}\textbf{\color{Purple}R5 · numeri sul colore}\\[6pt]
\includegraphics[width=\linewidth]{audit/evidenze/P07_F06_R5.pdf}
\end{minipage}\hfill
\begin{minipage}[t]{.48\textwidth}\textbf{\color{Teal}R6 · contrasto costante}\\[6pt]
\includegraphics[width=\linewidth]{figure/P07/F06.pdf}
\end{minipage}

\subsection*{Formalismo leggibile nella singola pagina}
In P11 sono state esplicitate le coordinate standardizzate, i centroidi, la norma e la trasposta impiegati nella derivazione del confine fra classi. In P10 la regola generale di tolleranza numerica è stata distinta dalla tolleranza esclusivamente assoluta che il codice applica alla rilettura del CSV. In M10 la condizione sulle radici dell’autoregressione è stata precisata rispetto alla causalità del processo.

Queste integrazioni permettono di collegare le espressioni matematiche alle operazioni illustrate, evitando che il lettore debba inferire il significato dei simboli dalla sola sequenza delle slide.

\subsection*{Registro linguistico}
La revisione conserva la prosa argomentata e interviene sui rinvii fra versioni, sulle formule introduttive generiche e sulle contrapposizioni stereotipate. Le negazioni che delimitano correttamente una conclusione statistica rimangono parte della spiegazione.

\newpage
\titolo{Verifiche e consultazione dei materiali}
\subsection*{Controlli della consegna R6}
\begin{tabular}{p{12.4cm}r}\toprule
Controllo & Esito\\\midrule
Fascicoli compilati, per complessive 861 pagine & 28/28\\
Pagine registrate nella rilettura individuale & 861/861\\
Interventi ricontrollati nel PDF & 357/357\\
Laboratori R: esito positivo e hash coerenti & 340/340\\
Estratti R delle slide: esito positivo e hash coerenti & 108/108\\
Laboratori Python introduttivi verificati & 20/20\\
Slide originali con raccordo documentato & 2150/2150\\\bottomrule
\end{tabular}

I controlli automatici verificano inoltre le dimensioni delle pagine, le altezze dei riquadri, l’assenza di righe eccedenti e caratteri mancanti, la disponibilità delle risorse richiamate e il contenimento delle etichette nei PDF delle figure. Il confronto dei file conferma che template originale e caratteri non sono stati modificati.

I resoconti di esecuzione preesistenti restano validi per i sorgenti il cui hash è invariato. Gli script R delle figure modificati e il laboratorio M12 N12 sono stati rieseguiti durante la revisione. I controlli e le versioni dell’ambiente sono inclusi nella cartella \texttt{verifiche}.

\subsection*{Che cosa è incluso}
La consegna comprende 261 figure R distinte, utilizzate in 262 posizioni; 232 diagrammi TikZ utilizzati nelle slide; 441 pagine con formule esplicite nei riquadri. Gli ulteriori simboli matematici nelle spiegazioni e nelle didascalie non sono inclusi in quest’ultimo conteggio.

\texttt{Indice\_corso.html} consente di aprire i fascicoli dopo l’estrazione dello ZIP. La cartella \texttt{audit} contiene questo rapporto, \texttt{Registro\_slide\_R6.html}, la versione CSV e il registro JSON. Il registro HTML permette di filtrare per modulo, cercare un argomento e visualizzare soltanto le pagine corrette.

\subsection*{Interpretazione dell’esito}
La verifica riguarda la lettura dei PDF e la loro coerenza didattica e grafica. Una prova con studenti non è stata svolta; la comprensione effettiva e la leggibilità da una specifica aula richiedono un riscontro nelle condizioni di utilizzo. Il registro rende comunque identificabile ogni pagina esaminata e ogni correzione apportata.

\end{document}
